\section{张量积}

\subsection{模的张量}

\begin{definition}[自由 Abel 群]\label{def:free-abel}
设集合 $S=\{s_i\}_{i\in I}$, 称集合
\[
  C(S)=\Bigl\{\textstyle\sum_{i\in I} m_i s_i \Bigm| m_i\in\mathbb{Z},\ s_i\in S,\ \text{只有有限项 } m_i\neq 0\Bigr\}
\]
在加法
\[
  \Bigl(\sum_{i\in I} m_i s_i\Bigr)+\Bigl(\sum_{i\in I} n_i s_i\Bigr)=\sum_{i\in I}(m_i+n_i)\,s_i
\]
下构成由 $S$ 生成的自由 Abel 群, 此时称 $S$ 为 $C(S)$ 的基.
\end{definition}

\begin{definition}[模的张量积]\label{def:tensor}
设 $R$ 是一个环, $M_R$ 为右 $R$-模, ${}_R N$ 为左 $R$-模.
记由集合 $M\times N$ 生成的自由 Abel 群为 $C(M\times N)$.
设 $L$ 为由以下三类元素生成的 $C(M\times N)$ 的子群:
\begin{enumerate}[label=(\arabic*)]
  \item $(m_1+m_2,n)-(m_1,n)-(m_2,n)$;
  \item $(m,n_1+n_2)-(m,n_1)-(m,n_2)$;
  \item $(m\cdot r,n)-(m,r\cdot n)$.
\end{enumerate}
其中 $r\in R$, $m,m_1,m_2\in M$, $n,n_1,n_2\in N$.
称商群 $C(M\times N)/L$ 为 $M_R$ 和 ${}_R N$ 在 $R$ 上的\textbf{张量积}, 记为 $M\otimes_R N$.
\end{definition}

\begin{remark}\label{rem:tensor-basic}
(i) $M\times N$ 仅仅是集合而不是模的直和, 我们没有在它上面定义运算, 故不存在"坐标分量对应相加"的说法; 只应将元素看作没有任何关系的形式记号.

(ii) $M\otimes_R N$ 为 Abel 群.

(iii) 有典范同态
\[
  \pi\colon C(M\times N)\longrightarrow M\otimes_R N,\qquad (m,n)\longmapsto \pi\bigl((m,n)\bigr),
\]
记 $\pi\bigl((m,n)\bigr)=m\otimes n$. 由 $\operatorname{Ker}(\pi)=L$ 可得
\[
  (m_1+m_2)\otimes n=m_1\otimes n+m_2\otimes n,\qquad
  m\otimes(n_1+n_2)=m\otimes n_1+m\otimes n_2,
\]
以及
\[
  (m\cdot r)\otimes n=m\otimes(r\cdot n).
\]
\end{remark}

\begin{remark}\label{rem:tensor-form}
$C(M\times N)$ 中的元素形如 $\sum_i b_i(m_i,n_i)$, $b_i\in\mathbb{Z}$, 故 $M\otimes_R N$ 中的元素形如
\[
  \pi\Bigl(\sum_i b_i(m_i,n_i)\Bigr)=\sum_i b_i\,(m_i\otimes n_i).
\]
对 $k\in\mathbb{Z}$ 有
\[
  (km)\otimes n=\underbrace{m\otimes n+\cdots+m\otimes n}_{|k|\ \text{个}}=k\,(m\otimes n),
\]
故元素可进一步写为 $\sum_i m_i\otimes n_i$ 的形式.
\end{remark}

\begin{example}\label{ex:tensor-basic}
(i) 考虑 $\mathbb{Z}\otimes_{\mathbb{Z}}\mathbb{Z}$. 注意
\[
  \sum_i m_i\otimes n_i=\sum_i 1\cdot m_i\otimes n_i=\sum_i 1\otimes m_i n_i=1\otimes\sum_i m_i n_i .
\]
因此设
\[
  \varphi\colon \mathbb{Z}\otimes_{\mathbb{Z}}\mathbb{Z}\longrightarrow\mathbb{Z},\qquad
  \sum_i m_i\otimes n_i=1\otimes\sum_i m_i n_i\longmapsto\sum_i m_i n_i .
\]
从 $\mathbb{Z}\times\mathbb{Z}$ 出发看 $\varphi$, 则 $(m,n)\mapsto mn$; 在 $C(\mathbb{Z}\times\mathbb{Z})$ 上为
$\sum_i b_i(m_i,n_i)\mapsto\sum_i b_i m_i n_i$.

下面验证 $\varphi$ 将 $L$ 中元素映为 $0$. 为方便, 将 $\varphi\bigl((m,n)\bigr)$ 简写为 $\varphi(m,n)$:
\begin{align*}
  \varphi\bigl((m_1+m_2,n)-(m_1,n)-(m_2,n)\bigr)
  &=(m_1+m_2)n-m_1n-m_2n=0,\\
  \varphi\bigl((m,n_1+n_2)-(m,n_1)-(m,n_2)\bigr)
  &=m(n_1+n_2)-mn_1-mn_2=0,\\
  \varphi\bigl((m\cdot r,n)-(m,r\cdot n)\bigr)&=0 .
\end{align*}
从而 $\varphi$ 是一个同态. 下面说明 $\varphi$ 是同构: 只需找 $\varphi'\colon\mathbb{Z}\to\mathbb{Z}\otimes_{\mathbb{Z}}\mathbb{Z}$
使 $\varphi'\varphi=\operatorname{Id}$. 定义
\[
  \varphi'\colon 1\longmapsto 1\otimes 1,\qquad a\longmapsto 1\otimes a=a\otimes 1,
\]
则
\[
  \varphi'\varphi\Bigl(\sum_i m_i\otimes n_i\Bigr)=\varphi'\Bigl(\sum_i m_i n_i\Bigr)
  =1\otimes\sum_i m_i n_i=\sum_i m_i\otimes n_i .
\]
于是 $\varphi'\varphi=\operatorname{Id}$, 故 $\mathbb{Z}\otimes_{\mathbb{Z}}\mathbb{Z}\cong\mathbb{Z}$.

(ii) $\mathbb{Z}\otimes_{\mathbb{Z}}\mathbb{Z}/n\mathbb{Z}\cong\mathbb{Z}/n\mathbb{Z}$. 定义
\[
  \varphi\colon \mathbb{Z}\otimes_{\mathbb{Z}}\mathbb{Z}/n\mathbb{Z}\longrightarrow\mathbb{Z}/n\mathbb{Z},\qquad
  \sum_i m_i\otimes[n_i]\longmapsto\sum_i m_i[n_i],
\]
\[
  \varphi'\colon \mathbb{Z}/n\mathbb{Z}\longrightarrow\mathbb{Z}\otimes_{\mathbb{Z}}\mathbb{Z}/n\mathbb{Z},\qquad
  [n]\longmapsto 1\otimes[n],
\]
即可.

(iii) 设 $R$ 为环, $M$ 为左 $R$-模, 则 $R\otimes_R M\cong M$. 定义
\[
  \varphi\colon R\otimes_R M\longrightarrow M,\qquad \sum_i r_i\otimes m_i\longmapsto\sum_i r_i\cdot m_i,
  \qquad
  \varphi'\colon M\longrightarrow R\otimes_R M,\qquad m\longmapsto 1\otimes m .
\]

(iv) 取右模 $\mathbb{Q}_{\mathbb{Z}}$ 和左模 ${}_{\mathbb{Z}}\mathbb{Z}/n\mathbb{Z}$, 则 $\mathbb{Q}\otimes_{\mathbb{Z}}\mathbb{Z}/n\mathbb{Z}=\{0\}$.
这是因为取 $\frac{p}{q}\in\mathbb{Q}$, $[s]\in\mathbb{Z}/n\mathbb{Z}$, 有
\[
  \frac{p}{q}\otimes[s]=\frac{p}{nq}\cdot n\otimes[s]=\frac{p}{nq}\otimes[ns]=0 .
\]

(v) $\mathbb{Z}/m\mathbb{Z}\otimes_{\mathbb{Z}}\mathbb{Z}/n\mathbb{Z}\cong\mathbb{Z}/\gcd(m,n)\mathbb{Z}$.
\end{example}

\subsection{泛性质}

\begin{definition}[平衡双加性映射]\label{def:balanced}
设 $R$ 是环, $M$ 为右 $R$-模, $N$ 为左 $R$-模, $L$ 为加群.
称 $f\colon M\times N\to L$ 为\textbf{平衡双加性映射}, 若
\begin{enumerate}[label=(\arabic*)]
  \item $f(m_1+m_2,n)=f(m_1,n)+f(m_2,n)$;
  \item $f(m,n_1+n_2)=f(m,n_1)+f(m,n_2)$;
  \item $f(m\cdot r,n)=f(m,r\cdot n)$.
\end{enumerate}
其中 $r\in R$, $m,m_1,m_2\in M$, $n,n_1,n_2\in N$.
\end{definition}

\begin{remark}\label{rem:balanced}
上述定义中, 双加性指 (i), (ii); 平衡指 (iii).
\end{remark}

\begin{theorem}[张量积的泛性质]\label{thm:tensor-universal}
设 $R$ 是环, $M_R$, ${}_R N$ 为模, $L$ 为加群.
对任意平衡双加性映射 $f\colon M\times N\to L$, 都存在唯一的加群同态 $\bar f\colon M\otimes_R N\to L$,
使得下图交换, 即 $\bar f\circ\pi=f$:
\[
  \begin{array}{ccc}
    M\times N & \xrightarrow{\ \pi\ } & M\otimes_R N\\[10pt]
    {\scriptstyle f}\Big\downarrow & \searrow{\ \scriptstyle\exists!\,\bar f} & \\[10pt]
    L & &
  \end{array}
\]
\end{theorem}

\begin{proof}
由平衡双加性映射 $f\colon M\times N\to L$ 可以找到加群同态
$\tilde f\colon C(M\times N)\to L$, 它由基 $M\times N$ 上的像确定, 即 $\tilde f(m,n)=f(m,n)$.
注意 $\tilde f$ 作用在张量积定义中的三类元素上有
\begin{align*}
  \tilde f\bigl((m_1+m_2,n)-(m_1,n)-(m_2,n)\bigr)
  &=f(m_1+m_2,n)-f(m_1,n)-f(m_2,n)=0,\\
  \tilde f\bigl((m,n_1+n_2)-(m,n_1)-(m,n_2)\bigr)
  &=f(m,n_1+n_2)-f(m,n_1)-f(m,n_2)=0,\\
  \tilde f\bigl((m\cdot r,n)-(m,r\cdot n)\bigr)
  &=f(m\cdot r,n)-f(m,r\cdot n)=0 .
\end{align*}
因此 $\tilde f$ 在 $L$ 上为零, 从而诱导出加群同态 $\bar f\colon M\otimes_R N\to L$.
事实上
\[
  \bar f\circ\pi\,(m,n)=\bar f(m\otimes n)=\tilde f(m,n)=f(m,n),
\]
即 $f=\bar f\circ\pi$. 又若 $g\colon M\otimes_R N\to L$ 满足 $g\circ\pi=f$, 则对任意 $m\otimes n$ 有
$g(m\otimes n)=f(m,n)=\bar f(m\otimes n)$, 故 $g=\bar f$, 唯一性成立.
\end{proof}

\begin{remark}\label{rem:tensor-universal}
(i) 典范映射
\[
  \pi\colon M\times N\longrightarrow M\otimes_R N,\qquad (m,n)\longmapsto m\otimes n
\]
本身就是平衡双加性映射.

(ii) 若 $L'$ 满足同样的泛性质, 则 $L'\cong M\otimes_R N$.
\end{remark}

\begin{remark}[证明的结构]\label{rem:tensor-proof-pic}
定理的证明可概括为下图:
\[
  \begin{array}{ccccc}
    M\times N & \xrightarrow{\ i\ } & C(M\times N) & \xrightarrow{\ \pi\ } & M\otimes_R N\\[10pt]
    & {\scriptstyle f}\searrow & {\scriptstyle \tilde f}\Big\downarrow & \swarrow{\ \scriptstyle \bar f} & \\[10pt]
    & & L & &
  \end{array}
\]
其中 $i$ 为嵌入, $\pi$ 为商映射, 上方复合为 $\pi\circ i$.
\end{remark}

\subsection{线性空间的张量积}

为引出线性空间的张量积, 先给出三个铺垫性的例子.

\begin{example}\label{ex:tensor-module-structure}
(i) 给定环 $R$, $S$, 取 $S$-$R$ 双模 ${}_S M_R$ 和左模 ${}_R N$.
此时 $M\otimes_R N$ 自然成为一个左 $S$-模:
\[
  S\times(M\otimes_R N)\longrightarrow M\otimes_R N,\qquad (s,m\otimes n)\longmapsto (s\cdot m)\otimes n .
\]
注意这里 $m\otimes n$ 是等价类, 故验证上述运算是否良定义是必要的.

(ii) 给定环 $R$, $S$, 取右模 $M_R$ 和 $R$-$S$ 双模 ${}_R N_S$.
类似地, $M\otimes_R N$ 是右 $S$-模.

(iii) 若 $R$ 是交换环, 则 $R=R^{\mathrm{op}}$, 此时 $M\otimes_R N$ 是 $R$-$R$ 双模.
\end{example}

\begin{remark}\label{rem:tensor-no-module}
一般而言, 给定右模 $M_R$ 和左模 ${}_R N$, 此时 $M\otimes_R N$ 仅仅是 Abel 群, 并没有模的结构.
\end{remark}

\begin{proposition}\label{prop:tensor-vector-space}
设 $K$ 是域, 给定线性空间 $M_K$ 和 ${}_K N$, 则 $M\otimes_K N$ 也是 $K$-线性空间.
\end{proposition}

\begin{definition}[双线性映射]\label{def:bilinear}
设 $K$ 是域, 给定线性空间 $M_K$, ${}_K N$, $V$.
称 $F\colon M\otimes_K N\to V$ 为\textbf{平衡双加性映射}, 若
\begin{enumerate}[label=(\arabic*)]
  \item $F(m_1+m_2,n)=F(m_1,n)+F(m_2,n)$;
  \item $F(m,n_1+n_2)=F(m,n_1)+F(m,n_2)$;
  \item $F(m\cdot k,n)=F(km,n)=F(m,kn)$.
\end{enumerate}
其中 $k\in K$, $m,m_1,m_2\in M$, $n,n_1,n_2\in N$.
\end{definition}

\begin{remark}\label{rem:bilinear}
若 $V$ 是线性空间, 则平衡性可加强为 $F(km,n)=kF(m,n)$;
此时平衡双加性映射即为\textbf{双线性映射}.
\end{remark}

\begin{theorem}[线性空间张量积的泛性质]\label{thm:tensor-universal-k}
设 $V_1,V_2$ 为 $K$-线性空间.
对任意双线性映射 $f\colon V_1\times V_2\to V_3$ ($V_3$ 为 $K$-线性空间),
都存在唯一的线性映射 $\bar f\colon V_1\otimes_K V_2\to V_3$, 使下图交换:
\[
  \begin{array}{ccc}
    V_1\times V_2 & \xrightarrow{\ \pi\ } & V_1\otimes_K V_2\\[10pt]
    {\scriptstyle f}\Big\downarrow & \searrow{\ \scriptstyle\exists!\,\bar f} & \\[10pt]
    V_3 & &
  \end{array}
\]
\end{theorem}

同构意义下, 给定维数的线性空间唯一. 为确定 $V_1\otimes_K V_2$ 的维数和基, 需要下面几个引理.

\begin{lemma}\label{lem:tensor-with-R}
设 $R$ 为环, $M$ 为左 $R$-模, 则作为左 $R$-模有 $R\otimes_R M\cong M$.
\end{lemma}

\begin{proof}
取
\[
  \varphi\colon R\otimes_R M\longrightarrow M,\qquad \sum_i r_i\otimes m_i\longmapsto\sum_i r_i m_i .
\]
对 $r\in R$ 有
\[
  \varphi\Bigl(r\cdot\sum_i r_i\otimes m_i\Bigr)=\varphi\Bigl(\sum_i rr_i\otimes m_i\Bigr)
  =\sum_i rr_i\,m_i=r\sum_i r_i m_i=r\cdot\varphi\Bigl(\sum_i r_i\otimes m_i\Bigr),
\]
故 $\varphi$ 是左 $R$-模同态, 从而是同构, 即 $R\otimes_R M\cong M$.
\end{proof}

\begin{corollary}\label{cor:tensor-KV}
设 $K$ 为域, $V$ 为 $K$-线性空间, 则有线性空间同构 $K\otimes_K V\cong V$.
\end{corollary}

\begin{lemma}\label{lem:tensor-direct-sum}
设 $R$ 为环, 取右模 $M_{1,R}$, $M_{2,R}$ 和左模 ${}_R N$, 则
\[
  (M_1\oplus M_2)\otimes_R N\cong (M_1\otimes_R N)\oplus(M_2\otimes_R N).
\]
\end{lemma}

\begin{proof}
令
\[
  \varphi\colon (M_1\oplus M_2)\otimes_R N\longrightarrow (M_1\otimes_R N)\oplus(M_2\otimes_R N),
  \qquad (m_1+m_2)\otimes n\longmapsto m_1\otimes n+m_2\otimes n,
\]
\[
  \psi\colon (M_1\otimes_R N)\oplus(M_2\otimes_R N)\longrightarrow (M_1\oplus M_2)\otimes_R N,
\]
\[
  m_1\otimes n_1+m_2\otimes n_2\longmapsto (m_1,n_1)\otimes n_1+(m_2,n_2)\otimes n_2,
\]
其中 $(m_1,m_2)\in M_1\oplus M_2$, $n,n_1,n_2\in N$. 直接验证 $\psi\varphi=\operatorname{Id}$ 且
$\varphi\psi=\operatorname{Id}$, 即可.
\end{proof}

\begin{remark}\label{rem:tensor-plus}
注意证明过程中"$+$"的意义: 第一个 $\oplus$ 中的加法是模的加法, 而
$(M_1\otimes_R N)\oplus(M_2\otimes_R N)$ 中的加法是分量相加.
\end{remark}

\begin{lemma}[结合律]\label{lem:tensor-assoc}
设 $R$, $S$ 为环, 取右模 $M_1$, 双模 ${}_R M_{2,S}$, 左模 ${}_S M_3$, 则
\[
  (M_1\otimes_R M_2)\otimes_S M_3=M_1\otimes_R(M_2\otimes_S M_3).
\]
\end{lemma}

\begin{proposition}\label{prop:tensor-dim}
设 $V_1,V_2$ 为 $K$-线性空间, 且 $\dim V_1=m$, $\dim V_2=n$, 那么
\begin{enumerate}[label=(\arabic*)]
  \item $\dim(V_1\otimes_K V_2)=mn$;
  \item 若 $V_1=\operatorname{span}_K\{\alpha_1,\dots,\alpha_m\}$,
        $V_2=\operatorname{span}_K\{\beta_1,\dots,\beta_n\}$, 则
        \[
          V_1\otimes_K V_2=\operatorname{span}_K\{\alpha_i\otimes\beta_j\mid 1\le i\le m,\ 1\le j\le n\}.
        \]
\end{enumerate}
\end{proposition}

\begin{proof}
任取 $\alpha=\sum_i k_i\alpha_i\in V_1$, $\beta=\sum_j l_j\beta_j\in V_2$, 则
$\alpha\otimes\beta\in V_1\otimes_K V_2$, 且有
\[
  \alpha\otimes\beta=\sum_{i,j} k_il_j\,\alpha_i\otimes\beta_j .
\]
另外, 记 $K\alpha_i=\operatorname{span}_K\{\alpha_i\}$, 则显然有 $V_1=\bigoplus_{i=1}^{m}K\alpha_i$.
于是由 Lemma \ref{lem:tensor-direct-sum} 与 Corollary \ref{cor:tensor-KV},
\[
  V_1\otimes_K V_2=\Bigl(\bigoplus_{i=1}^{m}K\alpha_i\Bigr)\otimes_K V_2
  \cong\bigoplus_{i=1}^{m}(K\alpha_i\otimes_K V_2)\cong\bigoplus_{i=1}^{m}V_2,
\]
故 $\dim(V_1\otimes_K V_2)=mn$, 且 $\{\alpha_i\otimes\beta_j\mid 1\le i\le m,\ 1\le j\le n\}$
为 $V_1\otimes_K V_2$ 的一组基.
\end{proof}
